\documentclass[10pt,a4paper]{amsart}
\usepackage{amsmath,amssymb,amsthm,graphicx}
\usepackage{verbatim}
\newenvironment{reference}{\comment}{\endcomment}
%\newenvironment{reference}{}{}
\newenvironment{slogan}{\comment}{\endcomment}
\newenvironment{history}{\comment}{\endcomment}

% For commutative diagrams we use Xy-pic
\usepackage[all]{xy}

% We use 2cell for 2-commutative diagrams.
\xyoption{2cell}
% Theorem environments.
%
\theoremstyle{plain}
\newtheorem{theorem}[subsection]{Theorem}
\newtheorem{proposition}[subsection]{Proposition}
\newtheorem{lemma}[subsection]{Lemma}

\theoremstyle{definition}
\newtheorem{definition}[subsection]{Definition}
\newtheorem{example}[subsection]{Example}
\newtheorem{exercise}[subsection]{Exercise}
\newtheorem{situation}[subsection]{Situation}

\theoremstyle{remark}
\newtheorem{remark}[subsection]{Remark}
\newtheorem{remarks}[subsection]{Remarks}

\numberwithin{equation}{subsection}

% Macros
%
\def\lim{\mathop{\mathrm{lim}}\nolimits}
\def\colim{\mathop{\mathrm{colim}}\nolimits}
\def\Spec{\mathop{\mathrm{Spec}}}
\def\Hom{\mathop{\mathrm{Hom}}\nolimits}
\def\Ext{\mathop{\mathrm{Ext}}\nolimits}
\def\SheafHom{\mathop{\mathcal{H}\!\mathit{om}}\nolimits}
\def\SheafExt{\mathop{\mathcal{E}\!\mathit{xt}}\nolimits}
\def\Sch{\mathit{Sch}}
\def\Mor{\mathop{\mathrm{Mor}}\nolimits}
\def\Ob{\mathop{\mathrm{Ob}}\nolimits}
\def\Sh{\mathop{\mathit{Sh}}\nolimits}
\def\NL{\mathop{N\!L}\nolimits}
\def\CH{\mathop{\mathrm{CH}}\nolimits}
\def\proetale{{pro\text{-}\acute{e}tale}}
\def\etale{{\acute{e}tale}}
\def\QCoh{\mathit{QCoh}}
\def\Ker{\mathop{\mathrm{Ker}}}
\def\Im{\mathop{\mathrm{Im}}}
\def\Coker{\mathop{\mathrm{Coker}}}
\def\Coim{\mathop{\mathrm{Coim}}}

% Boxtimes
%
\DeclareMathSymbol{\boxtimes}{\mathbin}{AMSa}{"02}

%
% Macros for moduli stacks/spaces
%
\def\QCohstack{\mathcal{QC}\!\mathit{oh}}
\def\Cohstack{\mathcal{C}\!\mathit{oh}}
\def\Spacesstack{\mathcal{S}\!\mathit{paces}}
\def\Quotfunctor{\mathrm{Quot}}
\def\Hilbfunctor{\mathrm{Hilb}}
\def\Curvesstack{\mathcal{C}\!\mathit{urves}}
\def\Polarizedstack{\mathcal{P}\!\mathit{olarized}}
\def\Complexesstack{\mathcal{C}\!\mathit{omplexes}}
% \Pic is the operator that assigns to X its picard group, usage \Pic(X)
% \Picardstack_{X/B} denotes the Picard stack of X over B
% \Picardfunctor_{X/B} denotes the Picard functor of X over B
\def\Pic{\mathop{\mathrm{Pic}}\nolimits}
\def\Picardstack{\mathcal{P}\!\mathit{ic}}
\def\Picardfunctor{\mathrm{Pic}}
\def\Deformationcategory{\mathcal{D}\!\mathit{ef}}
\usepackage{cleveref}
\crefname{originallemma}{original lemma}{original lemmas}
\Crefname{originallemma}{Original lemma}{Original lemmas}
\makeatletter
\newlabel{lemma-stalk-gives-point}{{04FM}{source}{}{sourcepointer-1}{}}
\newlabel{lemma-stalk-gives-point@cref}{{[originallemma][0][]04FM}{[1][0][]source}{}{}{}}
\makeatother
\crefname{originaldefinition}{original definition}{original definitions}
\Crefname{originaldefinition}{Original definition}{Original definitions}
\makeatletter
\newlabel{definition-stalk}{{040R}{source}{}{sourcepointer-2}{}}
\newlabel{definition-stalk@cref}{{[originaldefinition][0][]040R}{[1][0][]source}{}{}{}}
\makeatother
\crefname{originallemma}{original lemma}{original lemmas}
\Crefname{originallemma}{Original lemma}{Original lemmas}
\makeatletter
\newlabel{sites-lemma-point-site-topos}{{00YA}{source}{}{sourcepointer-3}{}}
\newlabel{sites-lemma-point-site-topos@cref}{{[originallemma][0][]00YA}{[1][0][]source}{}{}{}}
\makeatother
\crefname{originalproposition}{original proposition}{original propositions}
\Crefname{originalproposition}{Original proposition}{Original propositions}
\makeatletter
\newlabel{sites-proposition-point-limits}{{00YC}{source}{}{sourcepointer-4}{}}
\newlabel{sites-proposition-point-limits@cref}{{[originalproposition][0][]00YC}{[1][0][]source}{}{}{}}
\makeatother
\crefname{originalsection}{original section}{original sections}
\Crefname{originalsection}{Original section}{Original sections}
\makeatletter
\newlabel{sites-example-point-topological}{{00YE}{source}{}{sourcepointer-5}{}}
\newlabel{sites-example-point-topological@cref}{{[originalsection][0][]00YE}{[1][0][]source}{}{}{}}
\makeatother
\crefname{originalsection}{original section}{original sections}
\Crefname{originalsection}{Original section}{Original sections}
\makeatletter
\newlabel{sites-example-point-topology}{{04EJ}{source}{}{sourcepointer-6}{}}
\newlabel{sites-example-point-topology@cref}{{[originalsection][0][]04EJ}{[1][0][]source}{}{}{}}
\makeatother
\crefname{originalproposition}{original proposition}{original propositions}
\Crefname{originalproposition}{Original proposition}{Original propositions}
\makeatletter
\newlabel{proposition-etale-morphisms}{{03PC}{source}{}{sourcepointer-7}{}}
\newlabel{proposition-etale-morphisms@cref}{{[originalproposition][0][]03PC}{[1][0][]source}{}{}{}}
\makeatother
\crefname{originaldefinition}{original definition}{original definitions}
\Crefname{originaldefinition}{Original definition}{Original definitions}
\makeatletter
\newlabel{definition-geometric-point}{{03PO}{source}{}{sourcepointer-8}{}}
\newlabel{definition-geometric-point@cref}{{[originaldefinition][0][]03PO}{[1][0][]source}{}{}{}}
\makeatother
\crefname{originallemma}{original lemma}{original lemmas}
\Crefname{originallemma}{Original lemma}{Original lemmas}
\makeatletter
\newlabel{etale-lemma-etale-etale-local-technical}{{04HM}{source}{}{sourcepointer-9}{}}
\newlabel{etale-lemma-etale-etale-local-technical@cref}{{[originallemma][0][]04HM}{[1][0][]source}{}{}{}}
\makeatother
\crefname{originallemma}{original lemma}{original lemmas}
\Crefname{originallemma}{Original lemma}{Original lemmas}
\makeatletter
\newlabel{lemma-cofinal-etale}{{03PQ}{source}{}{sourcepointer-10}{}}
\newlabel{lemma-cofinal-etale@cref}{{[originallemma][0][]03PQ}{[1][0][]source}{}{}{}}
\makeatother
\crefname{originalsection}{original section}{original sections}
\Crefname{originalsection}{Original section}{Original sections}
\makeatletter
\newlabel{categories-section-codirected-limits}{{04AY}{source}{}{sourcepointer-11}{}}
\newlabel{categories-section-codirected-limits@cref}{{[originalsection][0][]04AY}{[1][0][]source}{}{}{}}
\makeatother
\crefname{originallemma}{original lemma}{original lemmas}
\Crefname{originallemma}{Original lemma}{Original lemmas}
\makeatletter
\newlabel{morphisms-lemma-value-at-one-point}{{04HB}{source}{}{sourcepointer-12}{}}
\newlabel{morphisms-lemma-value-at-one-point@cref}{{[originallemma][0][]04HB}{[1][0][]source}{}{}{}}
\makeatother
\title{Every point of the small etale site of a scheme comes from a geometric point}
\author{The Stacks Project Authors}
\date{stacks-project:etale-cohomology.tex:local-snapshot-20261011; GNU FDL 1.2 or later}
\begin{document}
\maketitle
\noindent\textit{Editorial note.} All delivered mathematical text is the author's own, copied from the pinned Stacks Project chapter file etale-cohomology.tex; no statement or proof step was rewritten, reordered or removed.\par\medskip
\noindent\textit{Editorial note.} The proof is the author's construction: the categorical properties, the comparison map, its well-definedness and the bijectivity are written out. Two routine steps are explicitly skipped in the counted span and are preserved verbatim there: the reduction of the construction from arbitrary objects to affine ones (source line 3725) and the verification of functoriality (source line 3787).\par\medskip
\noindent\textit{Editorial note.} Every labelled reference occurring inside the counted spans is carried in assembly.external\_references as a non-load-bearing navigation pointer (same-chapter ones at their definition site, cross-chapter ones at the invocation inside the counted span), so no load-bearing mathematics is replaced by a stub. No bibliography citation occurs inside the counted spans.\par\medskip
\noindent\textit{Editorial note.} The checked chapter list for same-name collisions is the set of stacks chapters staged in this repository: chow, cohomology, divisors, etale-cohomology, flat, more-algebra, more-morphisms, morphisms, spaces-more-morphisms, varieties together with this card's chapter etale-cohomology.\par\medskip
\noindent\textit{Editorial note.} The preamble is the chapter's own native preamble excerpt already used for this source: the Xy-pic setup (native preamble lines 14-18) and the native theorem declarations with the mathematical macros (native preamble lines 154-217).\par\medskip
\noindent\textit{Allowed standard prerequisites (editorial).}\begin{itemize}
\item Points of a site and points of its topos. Sites, Lemma sites-lemma-point-site-topos and Proposition sites-proposition-point-limits. Reduces to functors and supplies the properties of a point of a site.
\item The point attached to a geometric point. Same chapter, Lemma lemma-stalk-gives-point and its proof, and Definition definition-stalk. Supplies the reference point and the description of the underlying set used to construct the comparison map.
\item Generation of sheaves by representables. Topologies, Lemma topologies-lemma-alternative and Sites, Lemma sites-lemma-sheaf-coequalizer-representable. Supplies the final propagation to all sheaves.
\end{itemize}\medskip
\ifdefined\hypertarget\hypertarget{sourcepointer-1}{}\fi
\noindent\textit{Source reference.} Original lemma 04FM: sources/stacks-etale-cohomology/pinned-original-source.tex, source lines 3370--3370. Navigation pointer invoked at proof line 3588; the counted span contains the pointer only and the cited result is not reproduced inside it, so the pointer is registered as a non-load-bearing navigation stub in the established form for Stacks material. No bibliography citation is involved.\par\medskip
\ifdefined\hypertarget\hypertarget{sourcepointer-2}{}\fi
\noindent\textit{Source reference.} Original definition 040R: sources/stacks-etale-cohomology/pinned-original-source.tex, source lines 3341--3341. Navigation pointer invoked at proof line 3624; the counted span contains the pointer only and the cited result is not reproduced inside it, so the pointer is registered as a non-load-bearing navigation stub in the established form for Stacks material. No bibliography citation is involved.\par\medskip
\ifdefined\hypertarget\hypertarget{sourcepointer-3}{}\fi
\noindent\textit{Source reference.} Original lemma 00YA: sources/stacks-etale-cohomology/pinned-original-source.tex, source lines 3630--3630. Navigation pointer invoked at proof line 3630; the counted span contains the pointer only and the cited result is not reproduced inside it, so the pointer is registered as a non-load-bearing navigation stub in the established form for Stacks material. No bibliography citation is involved.\par\medskip
\ifdefined\hypertarget\hypertarget{sourcepointer-4}{}\fi
\noindent\textit{Source reference.} Original proposition 00YC: sources/stacks-etale-cohomology/pinned-original-source.tex, source lines 3634--3634. Navigation pointer invoked at proof line 3634; the counted span contains the pointer only and the cited result is not reproduced inside it, so the pointer is registered as a non-load-bearing navigation stub in the established form for Stacks material. No bibliography citation is involved.\par\medskip
\ifdefined\hypertarget\hypertarget{sourcepointer-5}{}\fi
\noindent\textit{Source reference.} Original section 00YE: sources/stacks-etale-cohomology/pinned-original-source.tex, source lines 3649--3649. Navigation pointer invoked at proof line 3649; the counted span contains the pointer only and the cited result is not reproduced inside it, so the pointer is registered as a non-load-bearing navigation stub in the established form for Stacks material. No bibliography citation is involved.\par\medskip
\ifdefined\hypertarget\hypertarget{sourcepointer-6}{}\fi
\noindent\textit{Source reference.} Original section 04EJ: sources/stacks-etale-cohomology/pinned-original-source.tex, source lines 3650--3650. Navigation pointer invoked at proof line 3650; the counted span contains the pointer only and the cited result is not reproduced inside it, so the pointer is registered as a non-load-bearing navigation stub in the established form for Stacks material. No bibliography citation is involved.\par\medskip
\ifdefined\hypertarget\hypertarget{sourcepointer-7}{}\fi
\noindent\textit{Source reference.} Original proposition 03PC: sources/stacks-etale-cohomology/pinned-original-source.tex, source lines 2592--2592. Navigation pointer invoked at proof line 3655; the counted span contains the pointer only and the cited result is not reproduced inside it, so the pointer is registered as a non-load-bearing navigation stub in the established form for Stacks material. No bibliography citation is involved.\par\medskip
\ifdefined\hypertarget\hypertarget{sourcepointer-8}{}\fi
\noindent\textit{Source reference.} Original definition 03PO: sources/stacks-etale-cohomology/pinned-original-source.tex, source lines 3201--3201. Navigation pointer invoked at proof line 3661; the counted span contains the pointer only and the cited result is not reproduced inside it, so the pointer is registered as a non-load-bearing navigation stub in the established form for Stacks material. No bibliography citation is involved.\par\medskip
\ifdefined\hypertarget\hypertarget{sourcepointer-9}{}\fi
\noindent\textit{Source reference.} Original lemma 04HM: sources/stacks-etale-cohomology/pinned-original-source.tex, source lines 3667--3667. Navigation pointer invoked at proof line 3667; the counted span contains the pointer only and the cited result is not reproduced inside it, so the pointer is registered as a non-load-bearing navigation stub in the established form for Stacks material. No bibliography citation is involved.\par\medskip
\ifdefined\hypertarget\hypertarget{sourcepointer-10}{}\fi
\noindent\textit{Source reference.} Original lemma 03PQ: sources/stacks-etale-cohomology/pinned-original-source.tex, source lines 3260--3260. Navigation pointer invoked at proof line 3706; the counted span contains the pointer only and the cited result is not reproduced inside it, so the pointer is registered as a non-load-bearing navigation stub in the established form for Stacks material. No bibliography citation is involved.\par\medskip
\ifdefined\hypertarget\hypertarget{sourcepointer-11}{}\fi
\noindent\textit{Source reference.} Original section 04AY: sources/stacks-etale-cohomology/pinned-original-source.tex, source lines 3708--3708. Navigation pointer invoked at proof line 3708; the counted span contains the pointer only and the cited result is not reproduced inside it, so the pointer is registered as a non-load-bearing navigation stub in the established form for Stacks material. No bibliography citation is involved.\par\medskip
\ifdefined\hypertarget\hypertarget{sourcepointer-12}{}\fi
\noindent\textit{Source reference.} Original lemma 04HB: sources/stacks-etale-cohomology/pinned-original-source.tex, source lines 3763--3763. Navigation pointer invoked at proof line 3763; the counted span contains the pointer only and the cited result is not reproduced inside it, so the pointer is registered as a non-load-bearing navigation stub in the established form for Stacks material. No bibliography citation is involved.\par\medskip
\begin{lemma}
\label{lemma-points-small-etale-site}
Let $S$ be a scheme.
\begin{enumerate}
\item Let $p$ be a point of the small \'etale site
$S_\etale$ of $S$ given by a functor
$u : S_\etale \to \textit{Sets}$.
Then there exists a geometric point $\overline{s}$ of $S$ such that
$p$ is isomorphic to the point of $S_\etale$ associated to
$\overline{s}$ in
Lemma \ref{lemma-stalk-gives-point}.
\item Let $p : \Sh(pt) \to \Sh(S_\etale)$ be a point
of the small \'etale topos of $S$. Then $p$ comes from a geometric point
of $S$, i.e., the stalk functor $\mathcal{F} \mapsto \mathcal{F}_p$
is isomorphic to a stalk functor as defined in
Definition \ref{definition-stalk}.
\end{enumerate}
\end{lemma}

\begin{proof}
By
Sites, Lemma \ref{sites-lemma-point-site-topos}
there is a one to one correspondence between points of the site and points
of the associated topos, hence it suffices to prove (1).
By
Sites, Proposition \ref{sites-proposition-point-limits}
the functor $u$ has the following properties:
(a) $u(S) = \{*\}$, (b) $u(U \times_V W) = u(U) \times_{u(V)} u(W)$, and
(c) if $\{U_i \to U\}$ is an \'etale covering, then
$\coprod u(U_i) \to u(U)$ is surjective.
In particular, if $U' \subset U$ is an open subscheme,
then $u(U') \subset u(U)$. Moreover, by
Sites, Lemma \ref{sites-lemma-point-site-topos}
we can write $u(U) = p^{-1}(h_U^\#)$, in other words $u(U)$ is the
stalk of the representable sheaf $h_U$. If
$U = V \amalg W$, then we see that $h_U = (h_V \amalg h_W)^\#$ and we get
$u(U) = u(V) \amalg u(W)$ since $p^{-1}$ is exact.

\medskip\noindent
Consider the restriction of $u$ to $S_{Zar}$. By
Sites, Examples \ref{sites-example-point-topological} and
\ref{sites-example-point-topology}
there exists a unique point $s \in S$ such that for $S' \subset S$ open we
have $u(S') = \{*\}$ if $s \in S'$ and $u(S') = \emptyset$ if $s \not \in S'$.
Note that if $\varphi : U \to S$ is an object of $S_\etale$ then
$\varphi(U) \subset S$ is open (see
Proposition \ref{proposition-etale-morphisms})
and $\{U \to \varphi(U)\}$ is an \'etale covering. Hence we conclude that
$u(U) = \emptyset \Leftrightarrow s \not \in \varphi(U)$.

\medskip\noindent
Pick a geometric point $\overline{s} : \overline{s} \to S$ lying over $s$, see
Definition \ref{definition-geometric-point}
for customary abuse of notation. Suppose that $\varphi : U \to S$ is an object
of $S_\etale$ with $U$ affine. Note that $\varphi$ is separated, and
that the fibre $U_s$ of $\varphi$ over $s$ is an affine scheme over
$\Spec(\kappa(s))$ which is the spectrum of a finite product of
finite separable extensions $k_i$ of $\kappa(s)$. Hence we may apply
\'Etale Morphisms, Lemma \ref{etale-lemma-etale-etale-local-technical}
to get an \'etale neighbourhood $(V, \overline{v})$ of $(S, \overline{s})$
such that
$$
U \times_S V = U_1 \amalg \ldots \amalg U_n \amalg W
$$
with $U_i \to V$ an isomorphism and $W$ having no point lying over
$\overline{v}$. Thus we conclude that
$$
u(U) \times u(V) =
u(U \times_S V) =
u(U_1) \amalg \ldots \amalg u(U_n) \amalg u(W)
$$
and of course also $u(U_i) = u(V)$. After shrinking $V$ a bit we can
assume that $V$ has exactly one point lying over $s$, and hence $W$ has no
point lying over $s$. By the above this then gives $u(W) = \emptyset$.
Hence we obtain
$$
u(U) \times u(V) =
u(U_1) \amalg \ldots \amalg u(U_n) =
\coprod\nolimits_{i = 1, \ldots, n} u(V)
$$
where the equality signs are compatible with the maps to the set $u(V)$.
Note that $u(V) \not = \emptyset$ as $s$ is in the image of $V \to S$.
In particular, we see that in this situation $u(U)$ is a finite
set with $n$ element, namely the left hand side has fibre $u(U)$ over
any element of $u(V)$ and the right hand side has fibre of cardinality $n$.

\medskip\noindent
Consider the limit
$$
\lim_{(V, \overline{v})} u(V)
$$
over the category of \'etale neighbourhoods $(V, \overline{v})$ of
$\overline{s}$. It is clear that we get the same value when taking
the limit over the subcategory of $(V, \overline{v})$ with $V$ affine.
By the previous paragraph (applied with the roles of $V$ and $U$ switched)
we see that in this case $u(V)$ is always a finite nonempty set.
Moreover, the limit is cofiltered, see
Lemma \ref{lemma-cofinal-etale}.
Hence by
Categories, Section \ref{categories-section-codirected-limits}
the limit is nonempty. Pick an element $x$ from this limit.
This means we obtain a $x_{V, \overline{v}} \in u(V)$ for
every \'etale neighbourhood $(V, \overline{v})$ of $(S, \overline{s})$
such that for every morphism of \'etale neighbourhoods
$\varphi : (V', \overline{v}') \to (V, \overline{v})$ we have
$u(\varphi)(x_{V', \overline{v}'}) = x_{V, \overline{v}}$.

\medskip\noindent
We will use the choice of $x$ to construct a functorial bijective map
$$
c : |U_{\overline{s}}| \longrightarrow u(U)
$$
for $U \in \Ob(S_\etale)$ which will conclude the proof. See
Lemma \ref{lemma-stalk-gives-point}
and its proof for a description of $|U_{\overline{s}}|$.
First we claim that it suffices to construct the map for $U$ affine.
We omit the proof of this claim.
Assume $U \to S$ in $S_\etale$ with $U$ affine, and let
$\overline{u} : \overline{s} \to U$ be an element of $|U_{\overline{s}}|$.
Choose a $(V, \overline{v})$ such that $U \times_S V$ decomposes
as in the third paragraph of the proof.
Then the pair $(\overline{u}, \overline{v})$ gives a geometric point of
$U \times_S V$ lying over $\overline{v}$ and determines one of the
components $U_i$ of $U \times_S V$. More precisely, there exists
a section $\sigma : V \to U \times_S V$ of the projection $\text{pr}_U$
such that $(\overline{u}, \overline{v}) = \sigma \circ \overline{v}$. Set
$c(\overline{u}) = u(\text{pr}_U)(u(\sigma)(x_{V, \overline{v}})) \in u(U)$.
We have to check this is independent of the choice of $(V, \overline{v})$. By
Lemma \ref{lemma-cofinal-etale}
the category of \'etale neighbourhoods is cofiltered.
Hence it suffice to show
that given a morphism of \'etale neighbourhood
$\varphi : (V', \overline{v}') \to (V, \overline{v})$ and a choice of a
section $\sigma' : V' \to U \times_S V'$ of the projection such that
$(\overline{u}, \overline{v'}) = \sigma' \circ \overline{v}'$
we have $u(\sigma')(x_{V', \overline{v}'}) = u(\sigma)(x_{V, \overline{v}})$.
Consider the diagram
$$
\xymatrix{
V' \ar[d]^{\sigma'} \ar[r]_\varphi & V \ar[d]^\sigma \\
U \times_S V' \ar[r]^{1 \times \varphi} &
U \times_S V
}
$$
Now, it may not be the case that this diagram commutes. The reason is
that the schemes $V'$ and $V$ may not be connected, and hence
the decompositions used to construct $\sigma'$ and $\sigma$ above may
not be unique. But we do know that
$\sigma \circ \varphi \circ \overline{v}' =
(1 \times \varphi) \circ \sigma' \circ \overline{v}'$
by construction. Hence, since $U \times_S V$ is \'etale over $S$,
there exists an open neighbourhood
$V'' \subset V'$ of $\overline{v'}$ such that the diagram does
commute when restricted to $V''$, see
Morphisms, Lemma \ref{morphisms-lemma-value-at-one-point}.
This means we may extend the diagram above to
$$
\xymatrix{
V'' \ar[r] \ar[d]^{\sigma'|_{V''}} &
V' \ar[d]^{\sigma'} \ar[r]_\varphi &
V \ar[d]^\sigma \\
U \times_S V'' \ar[r] &
U \times_S V' \ar[r]^{1 \times \varphi} &
U \times_S V
}
$$
such that the left square and the outer rectangle commute.
Since $u$ is a functor this implies that
$x_{V'', \overline{v}'}$ maps to the same element in
$u(U \times_S V)$ no matter which route we take through the
diagram. On the other hand, it maps to the elements
$x_{V', \overline{v}'}$ and $x_{V, \overline{v}}$ in
$u(V')$ and $u(V)$. This implies the desired equality
$u(\sigma')(x_{V', \overline{v}'}) = u(\sigma)(x_{V, \overline{v}})$.

\medskip\noindent
In a similar manner one proves that the construction
$c : |U_{\overline{s}}| \to u(U)$ is functorial in $U$;
details omitted. And finally, by the results of the
third paragraph it is clear that the map $c$ is bijective
which ends the proof of the lemma.
\end{proof}

\end{document}
